\begin{proposition}[2.1.8]
\label{IV.2.1.8-fr}
Sean $f:X\to Y$ un morfismo de preesquemas,
\[
0\to\mathscr F'\to\mathscr F\to\mathscr F''\to0
\]
una sucesión exacta de $\mathcal O_X$-módulos cuasicoherentes, tal que
$\mathscr F''$ sea plano sobre $Y$.
\begin{enumerate}
\item[\textup{(i)}] Para todo morfismo $g:Y'\to Y$ y todo $\mathcal O_{Y'}$-módulo
cuasicoherente $\mathscr G'$, la sucesión
$0\to\mathscr F'\otimes_Y\mathscr G'\to
\mathscr F\otimes_Y\mathscr G'\to
\mathscr F''\otimes_Y\mathscr G'\to0$
de $\mathcal O_{X'}$-módulos (donde $X'=X\times_Y Y'$) es exacta.
\item[\textup{(ii)}] Para que $\mathscr F$ sea plano sobre $Y$, es necesario y
suficiente que lo sea $\mathscr F'$.
\end{enumerate}
\end{proposition}

Evidentemente, se puede suponer que $X$, $Y$, $Y'$ son afines; la conclusión se
deduce entonces de (2.1.2) y de $(\mathbf{0}_{\mathrm I},6.1.2)$.

\begin{corollary}[2.1.9]
\label{IV.2.1.9-fr}
Sean $\mathscr L^\bullet$ un complejo de $\mathcal O_X$-módulos
cuasicoherentes, e $i$ un índice tal que, si se denota por
$d^i:\mathscr L^i\to\mathscr L^{i+1}$ el diferencial,
$\mathscr B^{i+1}(\mathscr L^\bullet)=\operatorname{Im}(d^i)$ y
$\mathscr Z'^{i+1}(\mathscr L^\bullet)=\operatorname{Coker}(d^i)$ sean
planos sobre $Y$. Entonces, con las notaciones de \textup{(2.1.8)}, el homomorfismo
canónico
\[
\mathscr H^i(\mathscr L^\bullet)\otimes_Y\mathscr G'
\longrightarrow
\mathscr H^i(\mathscr L^\bullet\otimes_Y\mathscr G')
\]
es biyectivo.
\end{corollary}

Como el producto tensorial es exacto por la derecha, se tiene
\[
\mathscr Z'^{i+1}(\mathscr L^\bullet)\otimes_Y\mathscr G'
=\operatorname{Coker}(d^i\otimes1)
=\mathscr Z'^{i+1}(\mathscr L^\bullet\otimes_Y\mathscr G')
\]
y
$\mathscr Z'^i(\mathscr L^\bullet)\otimes_Y\mathscr G'
=\mathscr Z'^i(\mathscr L^\bullet\otimes_Y\mathscr G')$. Además, en la
sucesión exacta
\[
0\to\mathscr B^{i+1}(\mathscr L^\bullet)\to\mathscr L^{i+1}
\to\mathscr Z'^{i+1}(\mathscr L^\bullet)\to0
\]
$\mathscr Z'^{i+1}(\mathscr L^\bullet)$ es plano sobre $Y$, por lo que de
(2.1.8, (i)) se deduce la sucesión exacta
\[
0\to\mathscr B^{i+1}(\mathscr L^\bullet)\otimes_Y\mathscr G'
\to\mathscr L^{i+1}\otimes_Y\mathscr G'
\to\mathscr Z'^{i+1}(\mathscr L^\bullet\otimes_Y\mathscr G')\to0
\]
y, por consiguiente,
$\mathscr B^{i+1}(\mathscr L^\bullet)\otimes_Y\mathscr G'
=\operatorname{Im}(d^i\otimes1)
=\mathscr B^{i+1}(\mathscr L^\bullet\otimes_Y\mathscr G')$. Entonces, puesto
que en la sucesión exacta
\[
0\to\mathscr H^i(\mathscr L^\bullet)\to\mathscr Z'^i(\mathscr L^\bullet)
\to\mathscr B^{i+1}(\mathscr L^\bullet)\to0
\]
$\mathscr B^{i+1}(\mathscr L^\bullet)$ es plano sobre $Y$, de (2.1.8, (i))
y de lo anterior se deduce la sucesión exacta
\[
0\to\mathscr H^i(\mathscr L^\bullet)\otimes_Y\mathscr G'
\to\mathscr Z'^i(\mathscr L^\bullet\otimes_Y\mathscr G')
\to\mathscr B^{i+1}(\mathscr L^\bullet\otimes_Y\mathscr G')\to0
\]
lo que demuestra el corolario.

\oldpage[IV-2]{8}
\begin{corollary}[2.1.10]
\label{IV.2.1.10-fr}
Sean $f:X\to Y$ un morfismo de preesquemas, $\mathscr F$ un
$\mathcal O_X$-módulo cuasicoherente y plano sobre $Y$,
$\mathscr L_\bullet=(\mathscr L_i)$ una resolución por la izquierda de
$\mathscr F$ formada por $\mathcal O_X$-módulos cuasicoherentes y planos sobre
$Y$. Entonces, para todo morfismo $g:Y'\to Y$ y todo
$\mathcal O_{Y'}$-módulo cuasicoherente $\mathscr G'$, el complejo
$\mathscr L_\bullet\otimes_Y\mathscr G'=
(\mathscr L_i\otimes_Y\mathscr G')_{i\geq0}$ es una resolución por la
izquierda de $\mathscr F\otimes_Y\mathscr G'$.

Además, si
$\mathscr Z_i(\mathscr L_\bullet)=\operatorname{Ker}
(\mathscr L_i\to\mathscr L_{i-1})$, los $\mathscr Z_i(\mathscr L_\bullet)$
son planos sobre $Y$, y se tiene
$\mathscr Z_i(\mathscr L_\bullet)\otimes_Y\mathscr G'
=\mathscr Z_i(\mathscr L_\bullet\otimes_Y\mathscr G')
=\operatorname{Ker}(\mathscr L_i\otimes_Y\mathscr G'
\to\mathscr L_{i-1}\otimes_Y\mathscr G')$.
\end{corollary}

Pongamos
$\mathscr R_i=\operatorname{Im}(\mathscr L_{i+1}\to\mathscr L_i)
=\mathscr Z_i(\mathscr L_\bullet)$; se tienen entonces las sucesiones exactas
\[
0\leftarrow\mathscr F\leftarrow\mathscr L_0\leftarrow\mathscr R_0\leftarrow0
\]
\[
\hbox to 12em{\dotfill}
\]
\[
0\leftarrow\mathscr R_i\leftarrow\mathscr L_{i+1}
\leftarrow\mathscr R_{i+1}\leftarrow0
\]
\[
\hbox to 12em{\dotfill}
\]
y, como $\mathscr F$ y los $\mathscr L_i$ son planos sobre $Y$, de
(2.1.8, (ii)) se deduce por recurrencia que todos los $\mathscr R_i$ son
también planos sobre $Y$; aplicando (2.1.8, (i)), se obtienen por tanto las
sucesiones exactas
\[
0\leftarrow\mathscr F\otimes_Y\mathscr G'
\leftarrow\mathscr L_0\otimes_Y\mathscr G'
\leftarrow\mathscr R_0\otimes_Y\mathscr G'\leftarrow0
\]
\[
0\leftarrow\mathscr R_i\otimes_Y\mathscr G'
\leftarrow\mathscr L_{i+1}\otimes_Y\mathscr G'
\leftarrow\mathscr R_{i+1}\otimes_Y\mathscr G'\leftarrow0
\qquad(i\geq0)
\]
que demuestran el corolario.

\begin{proposition}[2.1.11]
\label{IV.2.1.11-fr}
Sean $f:X\to Y$ un morfismo plano y $\mathscr F$ un $\mathcal O_Y$-módulo
cuasicoherente de presentación finita. Si $\mathscr J$ es el ideal de
$\mathcal O_Y$ que anula a $\mathscr F$, $f^*(\mathscr J)$ es el ideal de
$\mathcal O_X$ que anula a $f^*(\mathscr F)$.
\end{proposition}

En efecto, por definición se tiene una sucesión exacta $(\mathbf{0}_{\mathrm I},5.3.7)$
\[
0\to\mathscr J\to\mathcal O_Y\to
\mathscr H\!om_{\mathcal O_Y}(\mathscr F,\mathscr F)
\]
de donde, puesto que $f$ es plano, se obtiene una sucesión exacta
\[
0\to f^*(\mathscr J)\to\mathcal O_X\to
f^*(\mathscr H\!om_{\mathcal O_Y}(\mathscr F,\mathscr F))
\]
y, como por hipótesis $\mathscr F$ es un $\mathcal O_Y$-módulo de
presentación finita,
$f^*(\mathscr H\!om_{\mathcal O_Y}(\mathscr F,\mathscr F))$ se identifica
canónicamente con
$\mathscr H\!om_{\mathcal O_X}(f^*(\mathscr F),f^*(\mathscr F))$
$(\mathbf{0}_{\mathrm I},6.7.6)$, de donde se sigue la conclusión.

\begin{proposition}[2.1.12]
\label{IV.2.1.12-fr}
Sean $X$ un preesquema, $\mathscr F$ un $\mathcal O_X$-módulo de presentación
finita y $x$ un punto de $X$. Las condiciones siguientes son equivalentes:
\begin{enumerate}
\item[\emph{a)}] $\mathscr F_x$ es un $\mathcal O_x$-módulo plano.
\item[\emph{b)}] Existe una vecindad abierta $U$ de $x$ tal que
$\mathscr F|U$ sea un $(\mathcal O_X|U)$-módulo localmente libre.
\end{enumerate}
\end{proposition}

En efecto, $\mathscr F_x$ es un $\mathcal O_x$-módulo de presentación finita y
$\mathcal O_x$ es un anillo local; por tanto, es lo mismo decir que
$\mathscr F_x$ es un $\mathcal O_x$-módulo plano que decir que es un
$\mathcal O_x$-módulo libre (Bourbaki, \emph{Alg. comm.}, cap.~II, \S~3,
núm.~2, cor.~2 de la prop.~5); de ahí la conclusión, teniendo en cuenta
$(\mathbf{0}_{\mathrm I},5.2.7)$. Se observará que la proposición es válida para cualquier
\emph{espacio localmente anillado}.

\begin{proposition}[2.1.13]
\label{IV.2.1.13-fr}
Sea $f:X\to Y$ un morfismo de preesquemas. Si $f$ es plano en un punto
$x\in X$ y el anillo $\mathcal O_x$ es reducido (\textup{resp.} íntegro e integralmente
cerrado), entonces el anillo $\mathcal O_{f(x)}$
\oldpage[IV-2]{9}
es reducido (\textup{resp.} íntegro e integralmente cerrado). Si $f$ es fielmente plano
y $X$ es reducido (\textup{resp.} normal), entonces $Y$ es reducido (\textup{resp.} normal).
\end{proposition}

Pongamos $\mathcal O_{f(x)}=A$, $\mathcal O_x=B$. Si $B$ es un $A$-módulo
plano, también es un $A$-módulo fielmente plano $(\mathbf{0}_{\mathrm I},6.6.2)$, por lo
que $A$ se identifica con un subanillo de $B$; si $B$ es reducido, lo mismo
ocurre con $A$. Supongamos ahora que $B$ es íntegro e integralmente cerrado, y
sea $L$ su campo de fracciones; entonces el subanillo $A\subset B$ es íntegro;
denotemos por $K\subset L$ su campo de fracciones. La hipótesis implica que
$B\cap K=A$ (Bourbaki, \emph{Alg. comm.}, cap.~I, \S~3, núm.~5, prop.~10).
Si entonces $t\in K$ es entero sobre $A$, también lo es sobre $B$, por lo que
pertenece a $B$ por hipótesis y, por consiguiente, $t\in A$, lo cual demuestra
que $A$ es integralmente cerrado.

\begin{proposition}[2.1.14]
\label{IV.2.1.14-fr}
Sea $f:X\to Y$ un morfismo fielmente plano. Si $X$ es localmente íntegro y el
espacio subyacente a $Y$ es localmente noetheriano, entonces $Y$ es localmente
íntegro.
\end{proposition}

En efecto, todo $y\in Y$ es de la forma $f(x)$ para algún $x\in X$, y por
hipótesis $\mathcal O_y$ se identifica con un subanillo de $\mathcal O_x$
$(\mathbf{0}_{\mathrm I},6.6.1)$; como $\mathcal O_x$ es íntegro, lo mismo ocurre con
$\mathcal O_y$, y esto demuestra la proposición (\textbf I, 5.1.4).
